\section{Conclusiones}

Se estudió la atenuación de una señal debida al efecto Skin al hacerla atravesar una placa de aluminio. Se realizaron dos experimentos, en primer lugar se varió la frecuencia de la señal dejando el espesor fijo, y luego se cambió el espesor de la placa para una frecuencia fija. En ambos casos se obtuvieron valores de la conductividad eléctrica \(\sigma\) del aluminio y la profundidad de penetración \(\delta\).

En el primer experimento se obtuvo un valor de \(\sigma\) a partir de los datos del módulo de la función de transferencia de \(\sigma=8(1)\times 10^8 \) S/m, y a partir de los datos de la fase de \(\sigma=4(2)\times 10^7 \) S/m. El valor obtenido a partir de la fase es compatible con el valor tabulado para el aluminio de \(\sigma_{Al}=3\times 10^7 \) S/m, mientras que el valor obtenido a partir del módulo se encuentra fuera del orden de magnitud esperado. A pesar de esto, ambas series de datos siguen la tendencia predicha por el modelo teórico del efecto Skin. La discrepancia observada en el valor de \(\sigma\) obtenido a partir del módulo se puede atribuir a limitaciones del modelo teórico utilizado, que no contempla efectos geométricos e instrumentales presentes en el sistema experimental, siendo uno de estos el flujo magnético disperso que atraviesa regiones de aire entre la placa y las bobinas, o que, por la forma curva de las lineas de campo, llega a la bobina secundaria sin atravesar el material, lo que modifica la amplitud de la señal inducida.

En el segundo experimento, se obtuvo, a partir de un promedio ponderado de los resultados obtenidos en base a los datos del módulo y la fase, un valor de la profundidad de penetración de \(\delta=13.41(7)\) mm y una conductividad eléctrica de \(\sigma=4.26(3)\times 10^7 \) S/m. En el caso de la profundidad de penetración, el valor obtenido presenta una diferencia porcentual absoluta del \(20\%\) con el valor tabulado para el aluminio a \(50\) Hz, \(\delta_{Al}=11\) mm, mientras que el valor de la conductividad presenta una diferencia porcentual absoluta del \(13\%\) con la referencia. A pesar de estas diferencias, y de igual manera que en el primer experimento, se puede observar que la tendencia predicha por el modelo teórico del efecto Skin concuerda con los datos experimentales.